\section{Completing the approximation and its consequences}
\label{p14:reference:section}

The reconstruction is carried out in the algebra of signed measures.
For the parameters of the model its final reference is a probability.
The following positivity and continuity facts are elementary and do not
depend on the law approximation. The deduction of
Corollary~\ref{p1:cor:tv} from the main theorem is given below.

\subsection{Continuity of the reference}

For nonnegative $\mu,\mu'$,
\begin{equation}
 \|\Pi_\mu-\Pi_{\mu'}\|_1\le2|\mu-\mu'|.
 \label{p1:close:poisson-lipschitz}
\end{equation}
For example, if $\mu'=\mu+h$, couple by adding an independent
Poisson variable of mean $h$.  Its probability of being nonzero
is $1-e^{-h}\le h$, proving the inequality.
It follows that
\begin{equation}
 \|Q_{\mu,\alpha}-Q_{\mu',\alpha'}\|_1
 \le2(1+2|\alpha|)|\mu-\mu'|+2|\alpha-\alpha'|.
 \label{p1:close:reference-lipschitz}
\end{equation}
Here we have expanded the difference with coefficient $\alpha$
on $\Pi_\mu-\Pi_{\mu'}$ and used
$\|\Pi_{\mu'}*\Delta^{*2}\|_1\le4$.

\subsection{Positivity and the coefficient of total variation}
\label{p1:close:positivity-section}

The signed representation of the reference is useful in the
proof, but for the parameters of the theorem it is eventually
a probability.  The following exact criterion also specifies
its domain for general parameters.

\begin{lemma}[Positivity of the quadratic Poisson correction]
\label{p1:close:positivity}
Let $\mu>0$, $\alpha\in\mathbb R$, and
$\theta=\mu-\lfloor\mu\rfloor$.  Then $Q_{\mu,\alpha}$ is a
probability measure if and only if
\begin{equation}
 0\le\alpha\le
       \frac{2\mu^2}{\mu+\theta(1-\theta)}.
 \label{p1:close:positivity-range}
\end{equation}
Its masses are
\begin{equation}
 Q_{\mu,\alpha}\{k\}
 =e^{-\mu}\frac{\mu^k}{k!}
   \left[1+\frac{\alpha}{2\mu^2}
                    \bigl((k-\mu)^2-k\bigr)\right],
 \qquad k\ge0.
 \label{p1:close:charlier-masses}
\end{equation}
\end{lemma}

\begin{proof}
Formula \eqref{p1:close:charlier-masses} follows by expressing
the Poisson masses at $k-1$ and $k-2$ relative to the mass
at $k$; negative-index masses are taken to be zero.
The total mass is one because $\Delta^{*2}$ has total mass
zero.

If $\alpha<0$, the bracket is negative for all sufficiently
large $k$, so positivity is impossible.
For $\alpha\ge0$, minimize the quadratic
$f(k)=(k-\mu)^2-k$ over nonnegative integers.
Its successive difference is $f(k+1)-f(k)=2(k-\mu)$.
If $\mu$ is not an integer, the minimum is attained at
$k=\lceil\mu\rceil$; if $\mu$ is an integer, it is attained
at both $k=\mu$ and $k=\mu+1$.  In either case its value is
$-\mu-\theta(1-\theta)$.  Nonnegativity of the bracket
is therefore exactly \eqref{p1:close:positivity-range}.
\end{proof}

In the theorem, $\alpha_N>0$ for large $N$ and $\alpha_N\to0$,
while $\mu_N$ and $\lambda_N$ stay in a compact subinterval
of $(0,\infty)$.  Both corrected references are therefore
probabilities for all sufficiently large $N$.
The norm estimate simultaneously controls all their masses;
it is not only a pointwise assertion at each fixed count.

Define, for $\lambda>0$,
\begin{equation}
 \mathcal C(\lambda)
   =\frac14\|\Pi_\lambda*\Delta^{*2}\|_1
   =\frac14\E\left|
       1-\frac{2P_\lambda}{\lambda}
          +\frac{P_\lambda(P_\lambda-1)}{\lambda^2}
                   \right|,
 \qquad P_\lambda\sim\operatorname{Pois}(\lambda).
 \label{p1:close:charlier-constant}
\end{equation}
The coefficient is explicit even when $\lambda_N$ does not
converge. A summable signed sequence of total mass zero has
$\ell^1$ norm at least twice any positive coefficient. The mass
$e^{-\lambda}$ at zero of $\Pi_\lambda*\Delta^{*2}$ and its
convolution norm therefore give
\begin{equation}
 \frac{e^{-\lambda}}2\le\mathcal C(\lambda)\le1.
 \label{p1:close:constant-bounds}
\end{equation}

\subsection{Completion of the main theorem}
\label{p13:optimal:section}

The arithmetic and probabilistic estimates now meet at a
single quantity, the absolute third-order activity.
We complete the main theorem, then record what its precision
allows for unbounded scalar observables and for the choice
of mean.

\begin{proof}[Proof of Theorem~\ref{p1:thm:main}]
Write $\varepsilon_{\mathrm{out}}=\varepsilon$ for the requested
final margin and choose
$0<\tau<\min(\varepsilon_{\mathrm{out}}/4,1/960)$.
Invoke Theorem~\ref{p13:transfer:reconstruction} with accuracy $\tau$
and Corollary~\ref{p13:third:activity} with accuracy $\tau$.
The triple contribution is
\[
 N^\tau A_3\ll N^{-41/30+2\tau}
       =N^{-4/3-(1/30-2\tau)}.
\]
Since $1/30-2\tau>0$ and the other reconstruction term is
$N^{-4/3+\tau}$, this proves the upper bound.
Theorem~\ref{p13:lower:optimal} gives the lower bound for the same
exact parameters, using its own fixed Charlier interpolation order.

The mean estimate is Proposition~\ref{p1:arith:core}.
The pair asymptotic \eqref{p1:eq:alpha-leading} shows
that \(\alpha_N>0\), \(\alpha_N\to0\), and \(\mu_N\)
stays in a compact subset of \((0,\infty)\).
Lemma~\ref{p1:close:positivity} therefore makes the
reference a probability measure for sufficiently large \(N\).
Every auxiliary accuracy parameter and fixed order was chosen after
\(C,\varepsilon\) and before \(N\). The conclusion
is uniform on the entire fixed critical band.
\end{proof}

\subsection{Polynomial weights and fixed moments}

\begin{corollary}[Every fixed polynomial weight]
\label{p13:optimal:weighted}
For every fixed integer \(m\ge0\) and \(\varepsilon>0\),
\begin{equation}
 \sum_{k\ge0}(1+k)^m
   \left|\P(Z_N=k)-Q_{\mu_N,\alpha_N}\{k\}\right|
       \ll_{C,m,\varepsilon}N^{-4/3+\varepsilon}.
 \label{p13:optimal:weighted-bound}
\end{equation}
For every fixed \(r\ge2\), consequently,
\begin{equation}
 \E(Z_N)_r
 =\mu_N^r+\binom r2\mu_N^{r-2}\alpha_N
           +O_{C,r,\varepsilon}(N^{-4/3+\varepsilon}).
 \label{p13:optimal:moments}
\end{equation}
The equalities for \(r=0,1,2\) hold without a remainder
when the right side is interpreted in its natural form.
\end{corollary}

\begin{proof}
Let \(R=\Law Z_N-Q_{\mu_N,\alpha_N}\).
For every fixed integer \(v>1\), H\"older gives
\[
 \sum_k(1+k)^m|R\{k\}|
 \le\|R\|_1^{1-1/v}
       \left(\sum_k(1+k)^{mv}|R\{k\}|\right)^{1/v}.
\]
The second factor is uniformly bounded.
For the true law this follows from the independent
fixed-moment proof in Theorem~\ref{p1:fixed:moments};
for the reference it follows by dominating its variation
by a fixed sum of shifted Poisson laws, as in
Theorem~\ref{p13:lower:optimal}.
Choose \(v\) sufficiently large after \(\varepsilon\)
and then a sufficiently small margin in the main theorem.
This proves \eqref{p13:optimal:weighted-bound}.
The factorial generating function of \(Q_{\mu,\alpha}\)
is \(e^{\mu t}(1+\alpha t^2/2)\); its coefficient of
\(t^r\), multiplied by \(r!\), gives the moment formula.
\end{proof}

This interpolation concerns the scalar count and fixed polynomial degrees.

\begin{corollary}[No retained start]
\label{p13:optimal:zero}
Uniformly in a fixed critical band,
\[
 \P(Z_N=0)
   =e^{-\mu_N}\left(1+\frac{\alpha_N}{2}\right)
          +O_{C,\varepsilon}(N^{-4/3+\varepsilon}).
\]
\end{corollary}

\begin{proof}
Evaluate the approximation at zero, using
\((\Pi_\mu*\Delta^{*2})\{0\}=e^{-\mu}\).
\end{proof}

The event is the absence of a retained start in \(I_N\).
It does not exclude a long run that started before \(N\).

\subsection{Why the mean must be specified at this precision}

\begin{proposition}[The quantized mean bias]
\label{p13:optimal:nominal}
Let \(\eta_x\) indicate one-window compatibility and put
\(h_x=L-\dim S_x\). Then
\begin{equation}
 \mu_N-\lambda_N=pM_{N,L},\qquad
 M_{N,L}=\sum_{x\in I_N}(\eta_x2^{h_x}-1)\in\Z.
 \label{p13:optimal:integer-mean}
\end{equation}
Thus either \(\mu_N=\lambda_N\) or
\(|\mu_N-\lambda_N|\ge p\asymp_C N^{-1}\).
Keeping the same exact \(\alpha_N\), one has
\[
 \|\Law Z_N-Q_{\lambda_N,\alpha_N}\|_1
       \ll_{C,\varepsilon}N^{-1+\varepsilon}.
\]
On every sufficiently large admissible pair \((N,L)\)
for which \(M_{N,L}\ne0\),
\begin{equation}
 \|\Law Z_N-Q_{\lambda_N,\alpha_N}\|_1
       \asymp_C|\mu_N-\lambda_N|.
 \label{p13:optimal:nominal-distance}
\end{equation}
No assertion that \(M_{N,L}\ne0\) at infinitely many
scales is included.
\end{proposition}

\begin{proof}
Sum the one-window formula
\(\E J_x=\eta_xp2^{h_x}\) over the exactly \(N\)
centres to obtain \eqref{p13:optimal:integer-mean}.
For two means in the compact range in question,
\eqref{p1:close:reference-lipschitz} gives
\[
 \|Q_{\mu,\alpha_N}-Q_{\lambda,\alpha_N}\|_1
 \ll_C|\mu-\lambda|.
\]
The reverse bound follows from the mass at zero,
\(e^{-\theta}(1+\alpha_N/2)\), by the mean value theorem.
The upper bound for the nominal reference follows
from the main theorem and
\(|\mu_N-\lambda_N|\ll N^{-1+\varepsilon}\).

If \(M_{N,L}\ne0\), the mean difference is at least \(p\).
Choose a fixed margin less than \(1/3\) in the main
theorem; its remainder is then \(o(p)\).
The triangle and reverse triangle inequalities transfer
the two-sided reference comparison to
\eqref{p13:optimal:nominal-distance}.
\end{proof}

Irregular compatible and incompatible windows contribute
with opposite signs to \(M_{N,L}\). Their presence alone
does not prove a nonzero total bias.

\subsection{The total-variation coefficient}
\label{p14:comp:tv}

\begin{proof}[Proof of Corollary~\ref{p1:cor:tv}]
Theorem~\ref{p1:thm:main} and the reverse triangle
inequality first give
\[
 \TV(\Law Z_N,\Pi_{\mu_N})
 =\alpha_N\mathcal C(\mu_N)
            +O_{C,\varepsilon}(N^{-4/3+\varepsilon}).
\]
Use \(|\mu_N-\lambda_N|\ll N^{-1+\varepsilon}\) and
\eqref{p1:close:reference-lipschitz} to obtain
\[
 \Law Z_N-\Pi_{\lambda_N}
   =\frac{\alpha_N}{2}\Pi_{\lambda_N}*\Delta^{*2}
          +O_{\ell^1,C,\varepsilon}(N^{-1+\varepsilon}).
\]
The reverse triangle inequality, followed by the factor
$1/2$ converting the $\ell^1$ norm to total variation,
yields
\begin{equation}
 \TV(\Law Z_N,\Pi_{\lambda_N})
      =\alpha_N\mathcal C(\lambda_N)
          +O_{C,\varepsilon}(N^{-1+\varepsilon}).
 \label{p1:close:tv-exact-parameter}
\end{equation}
The pair asymptotic from
Theorem~\ref{p1:arith:pairs} now gives
\begin{equation}
 \TV(\Law Z_N,\Pi_{\lambda_N})
 \sim \mathcal C(\lambda_N)
       \frac{r_N+2}{18\,2^{r_N/3}}\,
       \lambda_N^{5/3}L N^{-2/3},
 \qquad r_N=L\bmod3.
 \label{p1:close:tv-asymptotic}
\end{equation}
The lower bound in \eqref{p1:close:constant-bounds} is uniform
on the fixed critical band, and the remainder is
$o(LN^{-2/3})$.  This proves the sharp order asserted in the
corollary.  The equivalence is relative to the displayed
oscillating expression; no convergence of its phase or
intensity factors is required.
\end{proof}

Replacing $\alpha_N$ by its leading equivalent proves
\eqref{p1:close:tv-asymptotic}.  It does not by itself preserve
the error $O(N^{-1+\varepsilon})$ in
\eqref{p1:close:tv-exact-parameter}, which uses the exact
second factorial cumulant.

Proposition~\ref{p1:discussion:charlier-norm} gives
$\|\Pi_1*\Delta^{*2}\|_1=3/e$, hence $\mathcal C(1)=3/(4e)$.
Along $N=2^L$ with $3\mid L$, the pair coefficient therefore gives
\begin{equation}
 \TV(\Law Z_N,\Pi_1)
       \sim \frac1{12e}\,L N^{-2/3}.
 \label{p1:close:dyadic-example}
\end{equation}
This is the coefficient in the example announced in the introduction.

The total-variation asymptotic uses the exact-marginal reconstruction
and the explicit pair coefficient. The independent positivity and
reference-continuity estimates do not use the reconstruction.
